\documentclass[../main.tex]{subfiles}

\begin{document}
\chapter{Riemannian Metrics}

In this chapter we officially define Riemannian metrics, and discuss some of the basic computational techniques associated with them.

\section{Definitions}
The Euclidean geometry is just $\mathbb{R}^n$ with an inner product $\langle \cdot, \cdot \rangle$. A general inner product on a vector space $V$ allows us to define lengths, angles, and orthogonality. The isomorphisms that preserve the inner product are called \textbf{linear isometries}.

\subsection{Riemannian Metrics}
To extend these geometric ideas to abstract smooth manifolds, we define a structure that amounts to a smoothly varying choice of inner product on each tangent space.

\begin{definition}{Riemannian Metric}{Riemannian Metric}
	Let $M$ be a smooth manifold. A \textbf{Riemannian metric} on $M$ is a smooth covariant 2-tensor field $g \in \mathcal{T}^2(M)$ whose value $g_p$ at each point $p \in M$ is an inner product on the tangent space $T_pM$. A smooth manifold equipped with a Riemannian metric is called a \textbf{Riemannian manifold}.
\end{definition}

\begin{proposition}{Existence of Riemannian Metrics}{Existence of Riemannian Metrics}
	Any smooth manifold admits a Riemannian metric.
\end{proposition}
\begin{proof}
	As a finite sum of inner products is an inner product, this follows from the existence of partitions of unity.
\end{proof}

This definition is also valid for smooth manifold with boundary. And many proofs of theorems in Riemannian geometry works for smooth manifolds with boundary as well. Many problems involving Riemannian manifolds with boundary can be addressed by embedding into a larger manifold without boundary and extending the Riemannian metric arbitrarily to the larger manifold.

\begin{remark}
	A Riemannian metric is not the same as a metric in the sense of metric spaces. For convenience, we shall just use the term \textbf{metric} to refer to a Riemannian metric.
\end{remark}

\begin{example}{Euclidean Metric}{Euclidean Metric}
	The Euclidean metric $\bar{g}$ is the Riemannian metric on $\mathbb{R}^n$ given by the standard inner product, under the natural identification of $T_p\mathbb{R}^n$ with $\mathbb{R}^n$ for each $p \in \mathbb{R}^n$. In coordinates, the Euclidean metric is given by
	\begin{equation*}
		\langle v, w \rangle_{\bar{g}} = \sum_{i=1}^n v^i w^i
	\end{equation*}
\end{example}
We shall always assume that we are working with Euclidean metric on $\mathbb{R}^n$ unless otherwise specified.

\subsection{Isometries}

We give the following definition of isometries between Riemannian manifolds.
\begin{definition}{Isometries}{Isometries}
	Suppose $(M, g)$ and $(\tilde{M}, \tilde{g})$ are Riemannian manifolds, with or without boundary. An \textbf{isometry} from $(M, g)$ to $(\tilde{M}, \tilde{g})$ is a diffeomorphism $\varphi: M \to \tilde{M}$ such that $\varphi^*\tilde{g} = g$. In other words, $\varphi$ is a smooth bijection such that $\mathrm{d} \varphi_p: T_pM \to T_{\varphi(p)}\tilde{M}$ is a linear isometry for each $p \in M$.

	We say that $(M, g)$ and $(\tilde{M}, \tilde{g})$ are \textbf{isometric} if there exists an isometry from $(M, g)$ to $(\tilde{M}, \tilde{g})$.
\end{definition}

A composition of isometries and the inverse of an isometry are again isometries, so being isometric is an equivalence relation on the class of Riemannian manifolds with or without boundary. Our subject, Riemannian geometry, is concerned primarily with properties of Riemannian manifolds that are preserved by isometries.

\begin{definition}{Local Isometries}{Local Isometries}
	Suppose $(M, g)$ and $(\tilde{M}, \tilde{g})$ are Riemannian manifolds, with or without boundary. A \textbf{local isometry} from $(M, g)$ to $(\tilde{M}, \tilde{g})$ is a smooth map $\varphi: M \to \tilde{M}$ such that for each $p \in M$, there exists an open neighborhood $U$ of $p$ such that $\varphi|_U: U \to \varphi(U)$ is an isometry onto its image.
\end{definition}

\begin{proposition}{Local Isometry Criterion}{Local Isometry Criterion}
	If $(M, g)$ and $(\tilde{M}, \tilde{g})$ are Riemannian manifolds of the same dimension, then a smooth map $\varphi: M \to \tilde{M}$ is a local isometry if and only if $\varphi^*\tilde{g} = g$.
\end{proposition}
\begin{proof}
	Easy, as $\mathrm{d} \varphi_p$ is a linear isometry and by the inverse function theorem, $\varphi$ is a local diffeomorphism.
\end{proof}

\begin{definition}{Flatness}{Flatness}
	A Riemannian $n$-manifold $(M, g)$ is said to be \textbf{flat} if for each $p \in M$, there exists an open neighborhood $U$ of $p$ and an isometry $\varphi: U \to V$ onto an open subset $V \subseteq \mathbb{R}^n$ with the Euclidean metric. We say that $(M, g)$ is \textbf{locally isometric} to $\mathbb{R}^n$ with the Euclidean metric.
\end{definition}

\begin{definition}{Isometry Group}{Isometry Group}
	Let $(M, g)$ be a Riemannian manifold. The \textbf{isometry group} of $(M, g)$ is the set of all isometries from $(M, g)$ to itself, with the operation of composition. We denote the isometry group by $\Iso(M, g)$.
\end{definition}

\begin{remark}
	Summer B. Myers and Norman E. Steenrod proved that if $M$ has finite many components, then $\Iso(M, g)$ have a finite-dimensional Lie group structure, acting smoothly on $M$.
\end{remark}

\subsection{Local Representation of Riemannian Metrics}
Let $(M, g)$ be a Riemannian manifold with or without boundary. Let $(x^1, \ldots, x^n)$ be a local coordinate system on $U \subseteq M$. Then the Riemannian metric $g$ can be expressed in terms of the local coordinates as
\begin{equation}
	g = g_{ij} \mathrm{d}x^i \otimes \mathrm{d}x^j
\end{equation}
for some smooth functions $g_{ij}: U \to \mathbb{R}$. The component functions of this tensor field constitute a matrix-valued function $(g_{ij})$ on $U$, with
\begin{equation*}
	g_{ij}(p) = \left\langle \frac{\partial}{\partial x^i}\Big|_p, \frac{\partial}{\partial x^j}\Big|_p \right\rangle_g
\end{equation*}
The matrix $(g_{ij})$ is symmetric and positive definite at each point $p \in U$, so we can express it as the symmetric product by
\begin{equation*}
	g = g_{ij} \mathrm{d}x^i \mathrm{d}x^j
\end{equation*}
For example, the Euclidean metric on $\mathbb{R}^n$ is given in the standard coordinates by
\begin{equation*}
	\bar{g} = \delta_{ij} \mathrm{d}x^i \mathrm{d}x^j
\end{equation*}

More generally, if $(E_1, \ldots, E_n)$ is a smooth local frame on $U$, and $(\epsilon^1, \ldots, \epsilon^n)$ is the dual coframe, then we can express the Riemannian metric as
\begin{equation}
	g = g_{ij} \epsilon^i \epsilon^j
\end{equation}

$g$ also acts on a pair of smooth vector fields $X, Y \in \mathfrak{X}(M)$ by
\begin{equation}
	g(X, Y) = g_{ij} X^i Y^j = \langle X, Y \rangle_g
\end{equation}
yeilding a smooth function on $M$.

A local frame $(E_1, \ldots, E_n)$ is said to be \textbf{orthonormal} if $g(E_i, E_j) = \delta_{ij}$ for all $i, j$. In this case, the Riemannian metric can be expressed as
\begin{equation*}
	g = (\epsilon^1)^2 + \cdots + (\epsilon^n)^2
\end{equation*}

\begin{proposition}{Existence of Orthonormal Frames}{Existence of Orthonormal Frames}
	Let $(M, g)$ be a Riemannian $n$-manifold with or without boundary. Then if $(X_j)$ is any smooth local frame for $TM$ over $U \subseteq M$, then there exists a smooth local frame $(E_j)$ on $U$ such that $\vspan\{E_1, \ldots, E_k\} = \vspan\{X_1, \ldots, X_k\}$ for each $k = 1, \ldots, n$, and $(E_j)$ is orthonormal with respect to $g$.

	In particular, every point $p \in M$ has a neighborhood $U$ on which there exists a smooth orthonormal frame $(E_1, \ldots, E_n)$ for $TM$.
\end{proposition}

\begin{remark}
	Warning: This does not say that we can find a smooth orthonormal \textbf{coordinate} frame. In fact, it is not possible to find a smooth orthonormal coordinate frame in general. Many would think that if we have a smooth orthonormal frame, then we can assign coordinates to the flow lines of the frame, and then we would have a smooth orthonormal coordinate frame. But this is not true in general because the flows may not commute. A concrete example would be the polar coordinates on $\mathbb{R}^2 \setminus \{0\}$, where the coordinate vector fields $\frac{\partial}{\partial r}$ and $\frac{\partial}{\partial \theta}$ do not commute.
\end{remark}

\begin{definition}{Unit Tangent Bundle}{Unit Tangent Bundle}
	For a Riemannian manifold $(M, g)$, the \textbf{unit tangent bundle} of $M$ is the set
	\begin{equation}
		UTM = \{(p, v) \in TM : g_p(v, v) = 1\}
	\end{equation}
\end{definition}

\begin{proposition}{Properties of the Unit Tangent Bundle}{Properties of the Unit Tangent Bundle}
	Let $(M, g)$ be a Riemannian manifold with or without boundary. Then $UTM$ is a smooth, properly embedded codimension-1 submanifold with boundary of $TM$. And $\partial(UTM) = \pi^{-1}(\partial M)$, where $\pi: TM \to M$ is the projection map.

	The unit tangent bundle is connected if and only if $M$ is connected if $n \geq 2$.

	The unit tangent bundle is compact if and only if $M$ is compact.
\end{proposition}
\begin{proof}
	We give the proof of each part.

	\textbf{codimension-1:} Use the level set theorem, as $F: TM \to \mathbb{R}$ defined by $F(p, v) = g_p(v, v)$ is smooth and $1$ is a regular value of $F$, then $UTM = F^{-1}(1)$ is a smooth codimension-1 submanifold of $TM$.

	\textbf{properly embedded:} A closed embedded submanifold is properly embedded: its inclusion into $TM$ is a proper map.

	\textbf{boundary:} Just use the boundary coordinate chart to check that $\partial(UTM) = \pi^{-1}(\partial M)$.

	Connectedness and compactness are easy to check.
\end{proof}

\section{Methods for Constructing Riemannian Metrics}
Many examples of Riemannian manifolds arise naturally as submanifolds, products, and quotients of other Riemannian manifolds. In this section, we introduce some of the tools for constructing such metrics.

\subsection{Riemannian Submanifolds}
Every submanifold of a Riemannian manifold automatically inherits a Riemannian metric, and many interesting Riemannian metrics are defined in this way.

\begin{lemma}{Riemannian Manifold by Immersion}{Riemannian Manifold by Immersion}
	Suppose $(\tilde{M}, \tilde{g})$ is a Riemannian manifold with or without boundary, and $M$ is a smooth manifold with or without boundary. If $F: M \to \tilde{M}$ is a smooth map, then the smooth covariant 2-tensor field $g = F^*\tilde{g}$ is a Riemannian metric on $M$ if and only if $F$ is an immersion.
\end{lemma}
\begin{proof}
	Simply from non degeneracy of
\end{proof}

This metric $g$ is said to be \textbf{induced} by the immersion $F$.And $F$ that satisfy $g = F^*\tilde{g}$ are called \textbf{isometric immersions}.

The most important examples of induced metrics occur on submanifolds. Suppose $M \subseteq \tilde{M}$ is an (immersed or embedded) submanifold of a Riemannian manifold $(\tilde{M}, \tilde{g})$. Then the inclusion map $\iota: M \hookrightarrow \tilde{M}$ is an immersion, and the induced metric $g = \iota^*\tilde{g}$ is a Riemannian metric on $M$. We say that $(M, g)$ is a \textbf{Riemannian submanifold} of $(\tilde{M}, \tilde{g})$. By
\begin{equation}
	g_p(v, w) = \tilde{g}_p(\mathrm{d} \iota_p(v), \mathrm{d} \iota_p(w))
\end{equation}
Because we usually identify $T_pM$ with a subspace of $T_p\tilde{M}$, we can write $\tilde{g}_p(v, w)$ instead of $\tilde{g}_p(\mathrm{d} \iota_p(v), \mathrm{d} \iota_p(w))$.

\begin{example}{Spheres}{Spheres}
	For $n \geq 1$, the $n$-sphere $S^n \subseteq \mathbb{R}^{n+1}$ is a smooth embedded submanifold of $\mathbb{R}^{n+1}$, and the Euclidean metric on $\mathbb{R}^{n+1}$ induces a Riemannian metric on $S^n$. This is the standard Riemannian metric on $S^n$, or called the \textbf{round metric}.
\end{example}

Now, if $(\tilde{M}, \tilde{g})$ is an $m$-dimensional smooth Riemannian manifold with or without boundary, and $M \subseteq \tilde{M}$ is an $n$-dimensional smooth submanifold with or without boundary, then a local frame $(\tilde{E}_1, \ldots, \tilde{E}_m)$ for $\tilde{M}$ over an open set $\tilde{U} \subseteq \tilde{M}$ is said to be \textbf{adapted} to $M$ if $(\tilde{E}_1, \ldots, \tilde{E}_n)$ is tangent to $M$. In case $\tilde{M}$ has empty boundary, the slice coordinates are available, adapted local orthonormal frames are easy to find.

\begin{proposition}{Existence of Adapted Orthonormal Frames}{Existence of Adapted Orthonormal Frames}
	Let $(\tilde{M}, \tilde{g})$ be an $m$-dimensional Riemannian manifold without boundary, and let $M \subseteq \tilde{M}$ be an $n$-dimensional embedded submanifold with or without boundary. Then for each point $p \in M$, there exists a neighborhood $\tilde{U} \subseteq \tilde{M}$ of $p$ and a smooth local orthonormal frame $(\tilde{E}_1, \ldots, \tilde{E}_m)$ on $\tilde{U}$ that is adapted to $M$.
\end{proposition}
\begin{proof}
	Apply the Gram–Schmidt algorithm to a coordinate frame in slice coordinates. Embedded manifolds imply the existence of slice coordinates.
\end{proof}

\begin{definition}{Normal}{Normal}
	Let $(\tilde{M}, \tilde{g})$ be a Riemannian manifold and let $M \subseteq \tilde{M}$ be a smooth submanifold with or without boundary. For any $p\in M$, a vector $v \in T_p\tilde{M}$ is said to be \textbf{normal} to $M$ at $p$ if $\forall w \in T_pM, \tilde{g}_p(v, w) = 0$.

	The \textbf{normal space} to $M$ at $p$ is the subspace of $T_p\tilde{M}$ consisting of all vectors normal to $M$ at $p$, denoted by $N_pM$. We can decompose
	\begin{equation}
		T_p\tilde{M} = T_pM \oplus N_pM
	\end{equation}
	The \textbf{normal bundle} of $M$ in $\tilde{M}$ is the vector bundle
	\begin{equation}
		NM = \bigsqcup_{p \in M} N_pM
	\end{equation}
\end{definition}

\begin{proposition}{The Normal Bundle}{The Normal Bundle}
	Let $(\tilde{M}, \tilde{g})$ be an $m$-dimensional Riemannian manifold and let $M \subseteq \tilde{M}$ be an $n$-dimensional immersed or embedded submanifold with or without boundary. Then the normal bundle $NM$ is a smooth rank-$(m-n)$ vector subbundle of $T\tilde{M}|_M$. There are smooth bundle homomorphisms
	\begin{equation}
		\pi^{\perp}: T\tilde{M}|_M \to NM, \quad \pi^{\top}: T\tilde{M}|_M \to TM
	\end{equation}
	called the \textbf{normal projection} and the \textbf{tangential projection}, respectively, such that for each $p \in M$, $\pi^{\perp}_p: T_p\tilde{M} \to N_pM$ and $\pi^{\top}_p: T_p\tilde{M} \to T_pM$ are the orthogonal projections onto $N_pM$ and $T_pM$, respectively.
\end{proposition}
\begin{proof}
	For any $p \in M$, there is a neighborhood $U \subseteq M$ of $p$ that is embedded in $\tilde{M}$. The rest follows from the existence of adapted orthonormal frames.
\end{proof}

\begin{remark}
	In case $\tilde{M}$ is a manifold with boundary, the preceding constructions do not always work, because there is not a fully general construction of slice coordinates in that case. However, there is a satisfactory result in case the submanifold is the boundary itself, using boundary coordinates in place of slice coordinates.
\end{remark}

\begin{proposition}{Existence of Outward-Pointing Normal}{Existence of Outward-Pointing Normal}
	If $(M, g)$ is a smooth Riemannian manifold with boundary, the normal bundle to $\partial M$ is a smooth rank-1 subbundle over $\partial M$, and there is a unique smooth outward-pointing unit normal vector field along all of $\partial M$.
\end{proposition}

Computations on a submanifold $M \subseteq \tilde{M}$ are carried out most conveniently in terms of a \textbf{smooth local parametrization} which is a smooth map $X: U \to \tilde{M}$ from an open set $U \subseteq \mathbb{R}^n$ such that $X(U)$ is an open subset of $M$ and $X: U \to X(U)$ is a diffeomorphism. Then for $V = X(U)$ and $\varphi = X^{-1}: V \to U$, we have $(V, \varphi)$ is a smooth local coordinate chart on $M$.

Suppose $(M, g)$ is a Riemannian submanifold of $(\tilde{M}, \tilde{g})$, and $X: U \to \tilde{M}$ is a smooth local parametrization of $M$. Then the coordinate representation of $g$ in these coordinates is given by
\begin{equation}
	(\varphi^{-1})^*g = X^*g = X^*\iota^*\tilde{g} = (\iota \circ X)^*\tilde{g}
\end{equation}

The simplicity of the formula for the pullback of a tensor field makes this expression exceedingly easy to compute, once a coordinate expression for $\tilde{g}$ is known. In particular, if $M$ is an immersed $n$-dimensional Riemannian submanifold of $\mathbb{R}^m$ and $X: U \to \mathbb{R}^n$ is a smooth local parametrization of $M$, then the induced metric on $U$ is given by
\begin{equation}
	g = X^*\bar{g} = \sum_{i=1}^m \mathrm{d}X^i \mathrm{d} X^i = \sum_{i=1}^m \left(\frac{\partial X^i}{\partial u^j} \mathrm{d}u^j\right)^2 = \sum_{i=1}^m \sum_{j, k=1}^n \frac{\partial X^i}{\partial u^j} \frac{\partial X^i}{\partial u^k} \mathrm{d}u^j \mathrm{d}u^k
\end{equation}

\begin{example}{Metrics in Graph Coordinates}{Metrics in Graph Coordinates}
	If $U \subseteq \mathbb{R}^n$ is open and $f: U \to \mathbb{R}$ is a smooth function, then the graph of $f$ is the subset
	\begin{equation}
		\Gamma(f) = \{(x, f(x)) : x \in U\} \subseteq \mathbb{R}^{n+1}
	\end{equation}
	which is an $n$-dimensional embedded submanifold of $\mathbb{R}^{n+1}$. It has a global parametrization $X: U \to \mathbb{R}^{n+1}$ given by $X(u) = (u, f(u))$. The corresponding coordinates $(u^1, \ldots, u^n)$ on $M = \Gamma(f)$ are called \textbf{graph coordinates}. The induced metric on $M$ is given in graph coordinates by
	\begin{equation*}
		X^*\bar{g} = (\mathrm{d} u^1)^2 + \cdots + (\mathrm{d}u^n)^2 + \mathrm{d} f^2
	\end{equation*}
\end{example}

\begin{example}{Surface of Revolution}{Surface of Revolution}
	Let $H$ be the half plane $\{(r, z) \in \mathbb{R}^2 : r > 0\}$, and suppose $C \subseteq H$ is an embedded 1-dimensional submanifold. The \textbf{surface of revolution} generated by $C$ is the subset
	\begin{equation}
		S_C = \left\{(x, y, z) \in \mathbb{R}^3 : \left(\sqrt{x^2 + y^2}, z\right) \in C\right\} \subseteq \mathbb{R}^3
	\end{equation}
	Here $C$ is called the \textbf{generating curve}. Every smooth local parametrization $\gamma(t) = (a(t), b(t))$ of $C$ gives rise to a smooth local parametrization for $S_C$ given by
	\begin{equation}
		X(t, \theta) = (a(t)\cos\theta, a(t)\sin\theta, b(t))
	\end{equation}
	at some small open set of $(t, \theta)$. The $t$-coordinate curves $t \mapsto X(t, \theta_0)$ are the \textbf{meridians}, and the $\theta$-coordinate curves $\theta \mapsto X(t_0, \theta)$ are the \textbf{parallels}. The induced metric on $S_C$ is given in these coordinates by
	\begin{equation*}
		X^*\bar{g} = (a'(t)^2 + b'(t)^2) \mathrm{d}t^2 + a(t)^2 \mathrm{d}\theta^2
	\end{equation*}
\end{example}

\begin{example}{$n$-Torus}{n-Torus}
	The smooth covering map $X: \mathbb{R}^n \to T^n$ restricts to a smooth local parametrization on any sufficiently small open subset of $\mathbb{R}^n$, and the induced metric is equal to the Euclidean metric in these coordinates. So the $n$-torus is flat.
\end{example}

\subsection{Riemannian Products}
Next, we discuss the construction of Riemannian metrics on products of Riemannian manifolds. Let $(M_1, g_1)$ and $(M_2, g_2)$ be Riemannian manifolds, the product manifold $M_1 \times M_2$ has a natural Riemannian metric $g = g_1 \oplus g_2$ called the \textbf{Riemannian product metric}, defined by
\begin{equation}
	g_{(p_1, p_2)}((v_1, v_2), (w_1, w_2)) = g_{1, p_1}(v_1, w_1) + g_{2, p_2}(v_2, w_2)
\end{equation}
In coordinates, if $(x^1, \ldots, x^{n})$ are local coordinates on $M_1$ and $(x^{n+1}, \ldots, x^{n+m})$ are local coordinates on $M_2$, then $(x^1, \ldots, x^{n+m})$ are local coordinates on $M_1 \times M_2$, and the Riemannian product metric is given by
\begin{equation}
	(g_{ij}) = \begin{pmatrix}
		(g_1)_{ij} & 0          \\
		0          & (g_2)_{ij}
	\end{pmatrix}
\end{equation}

Here is an important generalization of product metrics. Suppose $(M_1, g_1)$ and $(M_2, g_2)$ are Riemannian manifolds, and $f: M_1 \to \mathbb{R}_+$ is a \emph{strictly positive} smooth function. Then the \textbf{warped product} $M_1 \times_f M_2$ is the product manifold $M_1 \times M_2$ equipped with the Riemannian metric $g = g_1 \oplus f^2 g_2$, defined by
\begin{equation}
	g_{(p_1, p_2)}((v_1, v_2), (w_1, w_2)) = g_1|_{p_1}(v_1, w_1) + f(p_1)^2 g_2|_{p_2}(v_2, w_2)
\end{equation}

\begin{remark}
	Despite the name, a warped product is generally not a product metric unless $f$ is constant. The warped product metric can be visualized as smoothly scaling $M_2$ when attaching it to different points of $M_1$. A wide variety of metrics can be constructed in this way.
\end{remark}

\begin{example}{Warped Products}{Warped Products}
	\begin{itemize}
		\item With $f=1$, the warped product $M_1 \times_f M_2$ is just the Riemannian product $M_1 \times M_2$.
		\item Every surface of revolution can be expressed as a warped product. Let $H$ be the half plane $\{(r, z) \in \mathbb{R}^2 : r > 0\}$, and let $C \subseteq H$ be an embedded 1-dimensional submanifold. The surface of revolution $S_C$ generated by $C$ is diffeomorphic to the warped product $C \times_r S^1$, where $r: C \to \mathbb{R}_+$ is the projection onto the first coordinate. The diffeomorphism is given by $(p, e^{i\theta}) \mapsto (r(p)\cos\theta, r(p)\sin\theta, z(p))$. The induced metric on $S_C$ is equal to the warped product metric on $C \times_r S^1$.
		\item Let $\rho$ be the standard coordinate function on $\mathbb{R}_+$. Then the map $\Phi(\rho, \omega) = \rho \omega$ gives an isometry from the warped product $\mathbb{R}_+ \times_\rho S^{n-1}$ to $\mathbb{R}^n \setminus \{0\}$ with the Euclidean metric. This is the polar coordinates.
	\end{itemize}
\end{example}

\subsection{Riemannian Submersions}
Unlike submanifolds and products of Riemannian manifolds, which automatically inherit Riemannian metrics of their own, quotients of Riemannian manifolds inherit Riemannian metrics only under very special circumstances.

Suppose $\tilde{M}$ and $M$ are smooth manifolds, and $\pi: \tilde{M} \to M$ is a smooth submersion. Take $\tilde{g}$ to be a Riemannian metric on $\tilde{M}$. By the submersion level set theorem, each fiber $\tilde{M}_y = \pi^{-1}(y)$ is a properly embedded smooth submanifold of $\tilde{M}$, and at each point $x \in \tilde{M}$, we define two subspaces of the tangent space $T_x\tilde{M}$:
\begin{itemize}
	\item The \textbf{vertical subspace}
	      \begin{equation}
		      V_x = \ker(\mathrm{d}\pi_x) = T_x\tilde{M}_{\pi(x)}
	      \end{equation}
	\item The \textbf{horizontal subspace}
	      \begin{equation}
		      H_x = V_x^\perp
	      \end{equation}
\end{itemize}
Then $T_x\tilde{M} = V_x \oplus H_x$. Note that the vertical space is well defined for every submersion, because it does not refer to the metric; but the horizontal space depends on the metric.

A vector field on $\tilde{M}$ is said to be \textbf{vertical} if it is tangent to the fibers of $\pi$, and \textbf{horizontal} if it is everywhere tangent to the horizontal subspaces. Given a smooth vector field $X$ on $M$, a vector field $\tilde{X}$ on $\tilde{M}$ is said to be a \textbf{horizontal lift} of $X$ if $\tilde{X}$ is horizontal and $\pi$-related to $X$, i.e. $\mathrm{d}\pi_x(\tilde{X}_x) = X_{\pi(x)}$ for all $x \in \tilde{M}$.

The next proposition is the principal tool for doing computations on Riemannian submersions.

\begin{proposition}{Properties of Horizontal Vector Fields}{Properties of Horizontal Vector Fields}
	Let $\tilde{M}$ and $M$ be smooth manifolds, and let $\pi: \tilde{M} \to M$ be a smooth submersion. Let $\tilde{g}$ be a Riemannian metric on $\tilde{M}$.
	\begin{itemize}
		\item Every smooth vector field $W$ on $\tilde{M}$ can be uniquely decomposed as $W = W^H + W^V$, where $W^H$ is horizontal and $W^V$ is vertical.
		\item Every smooth vector field $X$ on $M$ has a unique horizontal lift $\tilde{X}$ to $\tilde{M}$.
		\item For any $x \in \tilde{M}$, and $v \in H_x$, there exists a vector field $X$ on $M$ whose horizontal lift $\tilde{X}$ satisfies $\tilde{X}_x = v$.
	\end{itemize}
\end{proposition}
\begin{proof}
	SORRY
\end{proof}

\begin{remark}
	It is important to be aware, though, that not every horizontal vector field on $\tilde{M}$ is the horizontal lift of a vector field on $M$. This is easy to understand geometrically: a horizontal lift should have the ``same'' vector field along each fiber, but a horizontal vector field on $\tilde{M}$ can vary along a fiber. As a counterexample, consider $\pi: \mathbb{R}^2 \to \mathbb{R}$ given by $\pi(x, y) = x$, and $W = y \partial_x$ on $\mathbb{R}^2$. Then $W$ is horizontal, but it is not the horizontal lift of any vector field on $\mathbb{R}$.
\end{remark}

Now we can identify some quotients of Riemannian manifolds that inherit metrics of their own.

\begin{definition}{Riemannian Submersions}{Riemannian Submersions}
	Suppose $(\tilde{M}, \tilde{g})$ and $(M, g)$ are Riemannian manifolds, and $\pi: \tilde{M} \to M$ is a smooth submersion. Then $\pi$ is said to be a \textbf{Riemannian submersion} if for each $x \in \tilde{M}$, the differential $\mathrm{d}\pi_x|_{H_x}: H_x \to T_{\pi(x)}M$ is a linear isometry. In other words, $\tilde{g}_x(v, w) = g_{\pi(x)}(\mathrm{d}\pi_x(v), \mathrm{d}\pi_x(w))$ for all $v, w \in H_x$.
\end{definition}
This just means that the metric $\tilde{g}$ ``behaves'' the same along each fiber of $\pi$, so that the metric $g$ on $M$ is well defined.

Given a Riemannian manifold $(\tilde{M}, \tilde{g})$ and a surjective submersion $\pi: \tilde{M} \to M$, it is almost never the case that there is a metric on $M$ that makes $\pi$ a Riemannian submersion, because this requires that the metric $\tilde{g}$ be ``consistent'' along each fiber of $\pi$. There is, however, an important special case in which there is such a metric. Suppose $\pi: \tilde{M} \to M$ is a smooth surjective submersion, and $G$ is a group acting on $\tilde{M}$. We say that the action is \textbf{vertical} if every element $\varphi \in G$ takes each fiber to itself, meaning that $\pi(\varphi \cdot x) = \pi(x)$ for all $x \in \tilde{M}$. We say that the action is \textbf{transitive on fibers} if for each $p, q \in \tilde{M}$ with $\pi(p) = \pi(q)$, there exists $\varphi \in G$ such that $\varphi \cdot p = q$.

If in addition $\tilde{M}$ has a Riemannian metric $\tilde{g}$, we say that the action is \textbf{isometric} if for each $\varphi \in G$, the map $x \mapsto \varphi \cdot x$ is an isometry. In that case, provided the action is effective (meaning that $\varphi \cdot x = x$ for all $x \in \tilde{M}$ implies $\varphi$ is the identity), we can identify $G$ with a subgroup of the isometry group $\Iso(\tilde{M}, \tilde{g})$. Since an isometry is, in particular, a diffeomorphism, every isometric action is an action by diffeomorphisms.

\begin{theorem}{Riemannian Submersions}{Riemannian Submersions}
	Let $(\tilde{M}, \tilde{g})$ be a Riemannian manifold, and let $\pi: \tilde{M} \to M$ be a smooth surjective submersion. Suppose $G$ is a group acting on $\tilde{M}$ that is vertical, transitive on fibers, and isometric. Then there exists a unique Riemannian metric $g$ on $M$ such that $\pi: (\tilde{M}, \tilde{g}) \to (M, g)$ is a Riemannian submersion.
\end{theorem}
\begin{proof}
	By intuition this fulfills the requirement that the metric $\tilde{g}$ be ``consistent'' along each fiber of $\pi$.
\end{proof}

\begin{remark}
	\textbf{The Role of the Group Action}

	The introduction of the group $G$ may at first seem unnecessary. Indeed, the essential issue is simply whether $\tilde g$ is consistent along the fibers of $\pi$, so that one can define
	\begin{equation*}
		g_p(u,v)=\tilde g_x(\tilde u,\tilde v),
	\end{equation*}
	where $x\in\pi^{-1}(p)$ and $\tilde u,\tilde v\in(\ker d\pi_x)^\perp$ are the horizontal lifts of $u,v\in T_pM$. The required condition is that this expression be independent of the choice of $x$ in the fiber. Thus, one could in principle assume this consistency directly.

	The purpose of the group action is to provide a natural mechanism for guaranteeing this condition. If $x,y$ lie in the same fiber, transitivity gives $\varphi\in G$ such that $y=\varphi\cdot x$. Since the action is vertical, $\pi\circ\varphi=\pi$, and hence
	\begin{equation*}
		d\pi_y\circ d\varphi_x=d\pi_x.
	\end{equation*}
	Thus $d\varphi_x$ carries horizontal lifts at $x$ to the corresponding horizontal lifts at $y$. Since the action is isometric,
	\begin{equation*}
		\tilde g_y(d\varphi_xX,d\varphi_xY)=\tilde g_x(X,Y),
	\end{equation*}
	so the horizontal inner products are the same at all points of a fiber.

	Therefore, the group action is not necessary for the existence of a Riemannian submersion; rather, it is a structural way to guarantee the necessary fiber-consistency. In many important examples, the group is already present because $M$ is constructed as a quotient $\tilde M/G$, in which case the fibers are the $G$-orbits and the consistency follows from the isometry of the action.
\end{remark}

The next corollary describes one important situation to which the preceding theorem applies.

\begin{corollary}{Lie Group Quotients}{Lie Group Quotients}
	Let $(\tilde{M}, \tilde{g})$ be a Riemannian manifold, and let $G$ be a Lie group acting smoothly, freely, properly, and isometrically on $\tilde{M}$. Then the orbit space $M = \tilde{M}/G$ has a unique smooth manifold structure and Riemannian metric such that $\pi: (\tilde{M}, \tilde{g}) \to (M, g)$ is a Riemannian submersion.
\end{corollary}
\begin{proof}
	The smooth structure follows the quotient manifold theorem, and the Riemannian metric follows from the preceding theorem.
\end{proof}

\begin{example}{The Fubini–Study Metric}{The FubiniStudy Metric}
	Let $n \in \mathbb{Z}_+$, and consider the complex projective space $\mathbb{CP}^n$. The map $\pi: \mathbb{C}^{n+1} \setminus \{0\} \to \mathbb{CP}^n$ by the equivalence relation $z \sim \lambda z$ for $\lambda \in \mathbb{C}^*$ is a smooth surjective submersion. Identifying $\mathbb{C}^{n+1}$ with $\mathbb{R}^{2n+2}$, we can equip $\mathbb{C}^{n+1} \setminus \{0\}$ with the Euclidean metric. Let $p$ be the restriction of $\pi$ to the unit sphere $S^{2n+1} \subseteq \mathbb{C}^{n+1}$. Then $p: S^{2n+1} \to \mathbb{CP}^n$ is a smooth surjective submersion, and the group $S^1$ acts smoothly, freely, properly, and isometrically on $S^{2n+1}$ by complex multiplication:
	\begin{equation*}
		\lambda \cdot (z^1, \ldots, z^{n+1}) = (\lambda z^1, \ldots, \lambda z^{n+1})
	\end{equation*}
	Therefore, by the preceding corollary, there exists a unique Riemannian metric on $\mathbb{CP}^n$ such that $p: S^{2n+1} \to \mathbb{CP}^n$ is a Riemannian submersion. This metric is called the \textbf{Fubini–Study metric}.
\end{example}

\subsection{Riemannian Coverings}
Another important special case of Riemannian submersions occurs in the context of covering maps.
\begin{definition}{Riemannian Coverings}{Riemannian Coverings}
	Suppose $(\tilde{M}, \tilde{g})$ and $(M, g)$ are Riemannian manifolds. A smooth covering map $\pi: \tilde{M} \to M$ is a \textbf{Riemannian covering} if $\pi$ is a local isometry.
\end{definition}

\begin{proposition}{Existence of Riemannian Coverings}{Existence of Riemannian Coverings}
	Suppose $\pi: \tilde{M} \to M$ is a smooth normal covering map (which means that the deck transformation group acts transitively on the fibers of $\pi$). And $\tilde{g}$ is any Riemannian metric on $\tilde{M}$ that is invariant under all deck transformations. Then there exists a unique Riemannian metric $g$ on $M$ such that $\pi: (\tilde{M}, \tilde{g}) \to (M, g)$ is a Riemannian covering.
\end{proposition}
\begin{proof}
	SORRY
\end{proof}

\begin{proposition}{Construct Normal Covering}{Construct Normal Covering}
	Suppose $(\tilde{M}, \tilde{g})$ is a Riemannian manifold, and $\Gamma$ is a \textbf{discrete Lie group} acting smoothly, freely, properly, and isometrically on $\tilde{M}$. Then the orbit space $M = \tilde{M}/\Gamma$ has a unique Riemannian metric $g$ such that the quotient map $\pi: \tilde{M} \to \tilde{M}/\Gamma$ is a normal Riemannian covering.
\end{proposition}

This gives
\begin{proposition}{}{}
	Suppose $(M, g)$ and $(\tilde{M}, \tilde{g})$ are connected Riemannian manifolds, and $\pi: \tilde{M} \to M$ is a normal Riemannian covering map, and $\Gamma = \Aut_\pi(\tilde{M})$ then $M$ is isometric to the quotient $\tilde{M}/\Gamma$.
\end{proposition}

\begin{example}{}{}
	The two element group $\Gamma = \{1, -1\}$ acts smoothly, freely, properly, and isometrically on the unit sphere $S^n$ by multiplication. The quotient space $S^n/\Gamma$ is the real projective space $\mathbb{RP}^n$, and the quotient map $q: S^n \to \mathbb{RP}^n$ is a normal Riemannian covering. As the action is isometric, there is a unique Riemannian metric on $\mathbb{RP}^n$ such that $q$ is a Riemannian covering.
\end{example}

\begin{example}{}{}
	The open Mobius band is a quotient $\mathbb{R}^2/\mathbb{Z}$ by $n \cdot (x, y) = (x+n, (-1)^n y)$, and the action is smooth, free, proper, and isometric. Therefore, the open Mobius band inherits a Riemannian metric from the Euclidean metric on $\mathbb{R}^2$.
\end{example}

\section{Basic Constructions on Riemannian Manifolds}
Throughout this section $M$ is a smooth manifold with or without boundary.

\subsection{Raising and Lowering Indices}
An inner product gives a natural isomorphism between a vector space and its dual by the Riesz representation theorem. In the context of Riemannian manifolds, the Riemannian metric $g$ does the same thing. We define a bundle homomorphism $\hat{g}: TM \to T^*M$ by
\begin{equation}
	\hat{g}(v)(w) = g_p(v, w)
\end{equation}
for all $v, w \in T_pM$. If $X,Y$ are smooth vector fields on $M$, then $\hat{g}(X)(Y) = g(X, Y)$. This means that $\hat{g}(X)$ is a smooth covector field (by the \textbf{Tensor Characterization Lemma}), and that $\hat{g}: \mathfrak{X}(M) \to \Omega^1(M)$ is a linear map over $C^\infty(M)$. This means that $\hat{g}$ is a smooth bundle homomorphism.

Given a smooth local frame $(E_1, \ldots, E_n)$ and its dual frame $(\epsilon^1, \ldots, \epsilon^n)$, let $g = g_{ij} \epsilon^i \epsilon^j$ be the local coordinate representation of $g$. If $X = X^i E_i$ is a smooth vector field, then
\begin{equation}
	\hat{g}(X) = g_{ij} X^i \epsilon^j = X_j \epsilon^j
\end{equation}
where we define
\begin{equation}
	X_j = g_{ij} X^i
\end{equation}
called the \textbf{lowered index} of $X$. With this in mind, the covector field $\hat{g}(X)$ is often denoted by $X^\flat$.

The inverse of the map $\hat{g}$ is denoted by $\hat{g}^{-1} = (g^{ij})$, and thus we have $g^{ik} g_{kj} = \delta^i_j$. In terms of a local frame, the inverse map is given by
\begin{equation}
	\hat{g}^{-1}(\omega) = \omega_i E^i, \qquad \omega_i = g^{ij} \omega_j
\end{equation}
called the \textbf{raised index} of $\omega$, denoted by $\omega^\sharp$. These two operations are called \textbf{lowering and raising indices}, or \textbf{musical isomorphisms}.

Probably the most important application of the sharp operator is to extend the classical gradient operator to Riemannian manifolds. If $g$ is a Riemannian metric on $M$, and $f \in C^\infty(M)$, then the \textbf{gradient} of $f$ is the smooth vector field
\begin{equation}
	\grad f = (\mathrm{d}f)^\sharp
\end{equation}
Unwinding the definitions, we see that
\begin{equation}
	\mathrm{d} f_p(w) = \langle \grad f |_p, w \rangle, \quad \forall p \in M, w \in T_pM
\end{equation}
and the local basis representation of $\grad f$ is given by
\begin{equation}
	\grad f = (g^{ij} E_i(f)) E_j
\end{equation}

The next proposition shows that the gradient has the same geometric interpretation on a Riemannian manifold as it does in Euclidean space. If $f$ is a smooth real-valued function on a smooth manifold $M$, recall $p\in M$ is called \textbf{a regular point} of $f$ if $\mathrm{d}f_p \neq 0$, and a \textbf{critical point} of $f$ if $\mathrm{d}f_p = 0$. And a level set $f^{-1}(c)$ is called \textbf{regular} if every point of $f^{-1}(c)$ is a regular point of $f$, and we have shown that each regular level set is an embedded smooth hypersurface of $M$.

\begin{proposition}{Gradient is Normal}{Gradient is Normal}
	Suppose $(M, g)$ is a Riemannian manifold, and $f \in C^\infty(M)$. Let $\mathcal{R} \subseteq M$ be the set of regular points of $f$. Then for each $c \in \mathbb{R}$, let the set $M_c = f^{-1}(c) \cap \mathcal{R}$. Then if $M_c \neq \emptyset$, it is an embedded smooth hypersurface of $M$, and the gradient $\grad f$ is everywhere normal to $M_c$.
\end{proposition}

The flat and sharp operators can be applied to tensors of any rank, in any index position, to convert tensors from covariant to contravariant or vice versa. We define them formally.

\begin{definition}{Musical Operation}{Musical Operation}
	If $F$ is any $(k, l)$-tensor field and $i \in \{1, \ldots, k+l\}$ is any covariant index position, we can form a new tensor field $F^\sharp$ of type $(k+1, l-1)$ by setting
	\begin{equation}
		F^\sharp(\alpha_1, \ldots, \alpha_{k+l}) = F(\alpha_1, \ldots, \alpha_{i-1}, \alpha_i^\sharp, \alpha_{i+1}, \ldots, \alpha_{k+l})
	\end{equation}
	whenever $\alpha_i$ are vector fields or covector fields as appropriate. Similarly we can define $F^\flat$ by
	\begin{equation}
		F^\flat(\alpha_1, \ldots, \alpha_{k+l}) = F(\alpha_1, \ldots, \alpha_{i-1}, \alpha_i^\flat, \alpha_{i+1}, \ldots, \alpha_{k+l})
	\end{equation}
\end{definition}
In coordinates, this is simply raising or lowering the $i$-th index of the tensor using the metric and its inverse.

Another important application of the flat and sharp operators is to extend the trace operator to covariant tensors. If $h$ is any covariant $k$-tensor field with $k \geq 2$, 2, we can raise one of its indices (say the last one for definiteness) and obtain a $(1, k-1)$-tensor field. We define the \textbf{trace} of $h$ to be
\begin{equation}
	\tr_g(h) = \tr(h^\sharp)
\end{equation}
The most important case is that of a covariant 2-tensor field, in which
\begin{equation}
	\tr_g(h) = g^{ij} h_{ij}
\end{equation}

\subsection{Inner Products of Tensors}
A Riemannian metric yields, by definition, an inner product on tangent vectors at each point. Because of the musical isomorphisms between vectors and covectors, it is easy to carry the inner product over to covectors as well.

Suppose $g$ is a Riemannian metric on $M$, and $x\in M$. We can define an inner product on the cotangent space $T^*_xM$ by
\begin{equation}
	\langle \omega, \eta \rangle_g = \langle \omega^\sharp, \eta^\sharp \rangle_g
\end{equation}
in coordinates
\begin{equation*}
	\langle \omega, \eta \rangle = g_{kl} (g^{ki} \omega_i) (g^{lj} \eta_j) = g^{ij} \omega_i \eta_j = \omega_i \eta^i = \omega^i \eta_i
\end{equation*}

This construction can be extended to tensor bundles of any rank, as the following proposition shows. If $E \to M$ is a smooth vector bundle, then a \textbf{smooth fiber metric} on $E$ is an inner product on each fiber $E_p$ that varies smoothly, in the sense that for any (local) smooth sections $\sigma, \tau$ of $E$, the inner product $\langle \sigma, \tau \rangle$ is a smooth function on $M$.

\begin{remark}
	Well, a smooth fiber metric is just a generalization of a Riemannian metric to arbitrary vector bundles.
\end{remark}

\begin{proposition}{Inner Products of Tensors}{Inner Products of Tensors}
	Let $(M, g)$ be an $n$-dimensional Riemannian manifold with or without boundary. There is a unique smooth fiber metric on each tensor bundle $T^{(k, l)}TM$ with the property that if $\alpha_1, \ldots , \alpha_{k+1}$ and $\beta_1, \ldots, \beta_{l+1}$ are smooth vector or covector fields as appropriate, then
	\begin{equation}
		\langle \alpha_1 \otimes \cdots \otimes \alpha_k, \beta_1 \otimes \cdots \otimes \beta_l \rangle_g = \prod_{i=1}^k \langle \alpha_i, \beta_i \rangle_g
	\end{equation}
	With this inner product, if $(E_1, \ldots, E_n)$ is a smooth local orthogonal frame for $TM$ and $(\epsilon^1, \ldots, \epsilon^n)$ is the dual frame for $T^*M$, then the collection of tensor fields $E_{i_1} \otimes \cdots \otimes E_{i_k} \otimes \epsilon^{j_1} \otimes \cdots \otimes \epsilon^{j_l}$ for $1 \leq i_1, \ldots, i_k, j_1, \ldots, j_l \leq n$ is a smooth local orthogonal frame for $T^{(k, l)}TM$. In terms of any (not necessarily orthonormal) frame, this fiber metric satisfies
	\begin{equation}
		\langle F, G \rangle = g_{i_1 r_1} \cdots g_{i_k r_k} g^{j_1 s_1} \cdots g^{j_l s_l} F^{i_1 \ldots i_k}_{j_1 \ldots j_l} G^{r_1 \ldots r_k}_{s_1 \ldots s_l}
	\end{equation}
\end{proposition}

\subsection{The Volume Form and Integration}
Another important construction provided by a metric on an oriented manifold is a canonical volume form.

\begin{proposition}{The Riemannian Volume Form}{The Riemannian Volume Form}
	Let $(M, g)$ be an oriented $n$-dimensional Riemannian manifold with or without boundary. Then there exists a unique smooth $n$-form $\mathrm{d} V_g$ on $M$, called the \textbf{Riemannian volume form}, such that:
	\begin{itemize}
		\item If $(\epsilon^1, \ldots, \epsilon^n)$ is any smooth local orthonormal coframe on $T^*M$, then
		      \begin{equation}
			      \mathrm{d} V_g = \epsilon^1 \wedge \cdots \wedge \epsilon^n
		      \end{equation}
		\item If $(E_1, \ldots, E_n)$ is any smooth local orthonormal frame on $TM$, then
		      \begin{equation}
			      \mathrm{d} V_g(E_1, \ldots, E_n) = 1
		      \end{equation}
		\item If $(x^1, \ldots, x^n)$ are local coordinates on $M$, then
		      \begin{equation}
			      \mathrm{d} V_g = \sqrt{\det(g_{ij})} \mathrm{d}x^1 \wedge \cdots \wedge \mathrm{d}x^n
		      \end{equation}
	\end{itemize}
\end{proposition}

The significance of the Riemannian volume form is that it allows us to integrate functions on an oriented Riemannian manifold, not just differential forms.

If $f$ is a continuous, compactly supported real-valued function on an oriented Riemannian $n$-manifold $(M, g)$ with or without boundary, then $f \mathrm{d} V_g$ is a compactly supported $n$-form, and we define the integral of $f$ over $M$ by
\begin{equation}
	\int_M f \mathrm{d} V_g
\end{equation}
If $M$ is compact, the volume of $M$ is defined to be
\begin{equation}
	\Vol_g(M) = \int_M \mathrm{d} V_g
\end{equation}

\begin{remark}
	Note the notation $\mathrm{d} V_g$ is used just by convenience, not meaning it is an exact form. In fact, if $M$ is a compact oriented Riemannian manifold without boundary, then $\mathrm{d} V_g$ is never exact. If it were exact, then by Stokes' theorem we would have integration of $\mathrm{d} V_g$ over $M$ equal to zero, which is impossible since the volume is positive.
\end{remark}

For Riemannian hypersurfaces, we have the following important characterization of the volume form on the hypersurface in terms of that of the ambient manifold. If $X$ is a vector field and $\mu$ is a differential form, recall $X \lrcorner \mu$ is the interior multiplication of $\mu$ by $X$, defined by $(X \lrcorner \mu)(X_1, \ldots, X_{k-1}) = \mu(X, X_1, \ldots, X_{k-1})$.

\begin{proposition}{Volume Form by Interior Multiplication}{Volume Form by Interior Multiplication}
	Suppose $M$ is a hypersurface in an oriented Riemannian manifold $(\tilde{M}, \tilde{g})$ and $g$ is the induced Riemannian metric on $M$. Then $M$ is orientable if and only if there exists a global unit normal vector field $N$ on $M$, and in that case, the Riemannian volume form $\mathrm{d} V_g$ on $M$ is given by
	\begin{equation}
		\mathrm{d} V_g = (N \lrcorner \mathrm{d} V_{\tilde{g}})|_M
	\end{equation}
\end{proposition}

When $M$ is not orientable, we can still define integrals of functions, but now we have to use densities instead of differential forms
\begin{proposition}{Riemannian Density}{Riemannian Density}
	If $(M, g)$ is a Riemannian manifold, then there exists a unique smooth positive density $\mu$ on $M$, called the \textbf{Riemannian density}, such that for ever local orthogonal frame $(E_1, \ldots, E_n)$ we have
	\begin{equation}
		\mu(E_1, \ldots, E_n) = 1
	\end{equation}
\end{proposition}
\begin{proof}
	$\mu$ can be defined in terms of any local orthogonal frame $(E_1, \ldots, E_n)$ by
	\begin{equation*}
		\mu = | \epsilon^1 \wedge \cdots \wedge \epsilon^n |
	\end{equation*}
\end{proof}

Let $(M, g)$ be a Riemannian manifold, with or without boundary. If $M$ is oriented and $\mathrm{d} V_g$ is the Riemannian volume form, then its Riemannian density is given by $\mu = |\mathrm{d} V_g|$. On the other hand, the Riemannian density is defined whether or not $M$ is orientable. It is customary to denote the Riemannian density by the same notation $\mathrm{d} V_g$ as the Riemannian volume form. In either case, we can define the integral of a compactly supported smooth function $f$ on $M$ by $\int_M f \mathrm{d} V_g$.

\subsection{The Divergence and Laplacian Operators}
Suppose $(M, g)$ is a oriented Riemannian $n$-manifold with or without boundary, and $\mathrm{d} V_g$ is the Riemannian volume form. If $X$ is a smooth vector field on $M$, then $X \lrcorner \mathrm{d} V_g$ is a smooth $(n-1)$-form on $M$. The exterior derivative of this $(n-1)$-form is a smooth $n$-form, so it can be expressed as a smooth function multiplied by the volume form $\mathrm{d} V_g$. This function is called the \textbf{divergence} of $X$, denoted by $\divergence X$, and is defined by
\begin{equation}
	\mathrm{d} (X \lrcorner \mathrm{d} V_g) = (\divergence X) \mathrm{d} V_g
\end{equation}

If $M$ is not orientable, we can still find a neighborhood $U$ of each point $p \in M$ that is orientable, and define $\divergence X$ on $U$ by the same formula. The result does not depend on the choice of orientation. In this way, we can define the divergence operator on any Riemannian manifold with or without boundary.

The most important application of the divergence operator is the \textbf{divergence theorem}.
\begin{theorem}{Divergence Theorem}{Divergence Theorem}
	Let $(M, g)$ be a compact Riemannian manifold with boundary, then for any smooth vector field $X$ on $M$, we have
	\begin{equation}
		\int_M \divergence X \mathrm{d} V_g = \int_{\partial M} \langle X, N \rangle \mathrm{d} V_{g|_{\partial M}}
	\end{equation}
	where $N$ is the outward unit normal vector field on $\partial M$.
\end{theorem}

The \textbf{Laplacian} is an operator on smooth functions $\Delta: C^\infty(M) \to C^\infty(M)$ defined by
\begin{equation}
	\Delta f = \divergence(\grad f)
\end{equation}

The next proposition gives alternative formulas for these operators.

\begin{proposition}{Coordinate Representation}{Coordinate Representation}
	Let $(M, g)$ be a Riemannian manifold with or without boundary, and let $(x^1, \ldots, x^n)$ be local coordinates on $M$. The coordinate representations of the divergence and Laplacian are as follows:
	\begin{equation}
		\divergence \left(X^i \frac{\partial}{\partial x^i}\right) = \frac{1}{\sqrt{\det(g)}} \frac{\partial}{\partial x^i} \left(\sqrt{\det(g)} X^i\right)
	\end{equation}
	\begin{equation}
		\Delta u = \frac{1}{\sqrt{\det(g)}} \frac{\partial}{\partial x^i} \left(\sqrt{\det(g)} g^{ij} \frac{\partial u}{\partial x^j}\right)
	\end{equation}
\end{proposition}


\section{Lengths and Distances}
Perhaps the most important tool that a Riemannian metric gives us is the ability to measure lengths of curves and distances between points. Throughout this section, $(M, g)$ is a Riemannian manifold with or without boundary.

Without further qualification, a \textbf{curve} in $M$ always means a parametrized curve, that is, a continuous map $\gamma: I \to M$ where $I \subseteq \mathbb{R}$ is some interval. A curve is said to be \textbf{smooth} if it is smooth as a map between smooth manifolds. If $I$ have endpoints, we say that $\gamma$ is smooth if it can be extended to a smooth map on an open interval containing $I$. A \textbf{curve segment} is a curve whose domain is a compact interval.

A smooth curve $\gamma: I \to M$ has a well-defined velocity $\gamma'(t) \in T_{\gamma(t)}M$ for each $t \in I$. We say that $\gamma$ is \textbf{regular} if $\gamma'(t) \neq 0$ for all $t \in I$, which means $\gamma$ is an immersion, so its image has no corners or kinks. We wish to use curve segments as measuring sticks to define distances between points in a connected Riemannian manifold. Many aspects of the theory become technically much simpler if we work with a slightly larger class of curve segments instead of just the regular ones. We now describe the appropriate class of curves.

If $[a, b] \subseteq \mathbb{R}$ is a compact interval, a partition of $[a, b]$ is a finite sequence of points $a = a_0 < a_1 < \cdots < a_k = b$. If $M$ is a smooth manifold with or without boundary, a (continuous) curve segment $\gamma: [a, b] \to M$ is said to be \textbf{piecewise smooth} if there exists a partition $a = a_0 < a_1 < \cdots < a_k = b$ such that $\gamma|_{[a_{i-1}, a_i]}$ is a regular smooth curve segment for each $i = 1, \ldots, k$. For brevity, we refer to a piecewise regular curve segment as an \textbf{admissible curve}, and any partition of $[a, b]$ that makes $\gamma$ piecewise smooth is called an \textbf{admissible partition} of $\gamma$. For each admissible partition, at each intermediate point $a_i$, there are two \textbf{one-sided velocities}
\begin{equation}
	\gamma'(a_i^-) = \lim_{t \nearrow a_i} \gamma'(t), \qquad \gamma'(a_i^+) = \lim_{t \searrow a_i} \gamma'(t)
\end{equation}
They are both nonzero, but they need not be equal.
\begin{remark}
	There are many admissible partitions for a given admissible curve, because we can always add more points to the partition. It is also convenient to consider any map $\gamma: \{a\} \to M$ whose domain is a single real number to be an admissible curve.
\end{remark}

If $\gamma: I \to M$ is a smooth curve, we define a \textbf{reparameterization} of $\gamma$ to be a smooth curve $\tilde{\gamma} = \gamma \circ \phi: I' \to M$ where $\phi: I' \to I$ is a diffeomorphism. Because intervals are connected, $\varphi$ is either strictly increasing or strictly decreasing. If $\varphi$ is strictly increasing, we say that $\tilde{\gamma}$ is a \textbf{forward reparameterization} of $\gamma$, and if $\varphi$ is strictly decreasing, we say that $\tilde{\gamma}$ is a \textbf{backward reparameterization} of $\gamma$. For an admissible curve, we say that $\tilde{\gamma}$ is a reparameterization of $\gamma$ if it parts into reparameterizations of each smooth segment of $\gamma$.

\begin{definition}{Length of Curve}{Length of Curve}
	If $\gamma: [a, b] \to M$ is an admissible curve, we define its \textbf{length} to be
	\begin{equation}
		L_g(\gamma) = \int_a^b |\gamma'(t)|_g \mathrm{d}t
	\end{equation}
\end{definition}
The integrand is bounded and continuous everywhere on $[a, b]$ except at the finitely many points so this integral is well-defined.

\begin{proposition}{Properties of Lengths}{Properties of Lengths}
	Suppose $(M, g)$ is a Riemannian manifold with or without boundary, and $\gamma: [a, b] \to M$ is an admissible curve.
	\begin{itemize}
		\item \textbf{Additivity:} If $a < c < b$, then
		  \begin{equation}
			  L_g(\gamma) = L_g(\gamma|_{[a, c]}) + L_g(\gamma|_{[c, b]})
		  \end{equation}
		\item \textbf{Reparameterization Invariance:} If $\tilde{\gamma}$ is a reparameterization of $\gamma$, then
		  \begin{equation}
			  L_g(\tilde{\gamma}) = L_g(\gamma)
		  \end{equation}
		\item \textbf{Isometry Invariance:} If $\varphi: M \to \tilde{M}$ is a local isometry between Riemannian manifolds $(M, g)$ and $(\tilde{M}, \tilde{g})$, then
		  \begin{equation}
		  	L_g(\gamma) = L_{\tilde{g}}(\varphi \circ \gamma)
		  \end{equation}
	\end{itemize}
\end{proposition}

The \textbf{arc-length funcion} of an admissible curve $\gamma: [a, b] \to M$ is the function $s: [a, b] \to \mathbb{R}$ defined by
\begin{equation}
	s(t) = L_g(\gamma|_{[a, t]}) = \int_a^t |\gamma'(\tau)|_g \mathrm{d}\tau
\end{equation}
It is continuous everywhere, and it follows from the fundamental theorem of calculus that it is smooth wherever $\gamma$ is smooth, with $s'(t) = |\gamma'(t)|_g$.

For this reason, if $\gamma$ is any smooth curve (not necessarily a curve segment), we define the \textbf{speed} of $\gamma$ at $t \in I$ to be $|\gamma'(t)|_g$. We say that $\gamma$ is a \textbf{unit-speed curve} if $|\gamma'(t)|_g = 1$ for all $t \in I$, and a \textbf{constant-speed curve} if $|\gamma'(t)|_g$ is constant for all $t \in I$. The definition generalized to piecewise smooth curves by applying it to each smooth segment.

If $\gamma: [a, b] \to M$ is a unit-speed admissible curve, then its arc-length function has the simple form $s(t) = t - a$. If in addition $a=0$, then $s(t) = t$ and then for this reason, a unit-speed admissible curve whose parameter interval is of the form $[0, b]$ is often called to be \textbf{parameterized by arc length}.

\begin{proposition}{Unit Speed Reparameterization}{Unit Speed Reparameterization}
	Suppose $(M, g)$ is a Riemannian manifold with or without boundary,
	\begin{itemize}
		\item Every regular curve in $M$ has a unit-speed forward reparameterization.
		\item Every admissible curve in $M$ has a unique forward reparameterization by arc length.
	\end{itemize}
\end{proposition}

\subsection{The Riemannian Distance Function}
We now define the Riemannian distance function on a connected Riemannian manifold.

\begin{definition}{Riemannian Distance Function}{Riemannian Distance Function}
	Suppose $(M, g)$ is a connected Riemannian manifold with or without boundary. For any two points $p, q \in M$, we define the \textbf{Riemannian distance} between $p$ and $q$ to be
	\begin{equation}
		d_g(p, q) = \inf \{ L_g(\gamma) : \gamma \text{ is an admissible curve from } p \text{ to } q \}
	\end{equation}
\end{definition}
The following proposition guarantees that $d_g$ is well-defined.

\begin{proposition}{Existence of Admissible Curves}{Existence of Admissible Curves}
	Suppose $M$ is a connected smooth manifold with or without boundary. Then for any two points $p, q \in M$, there exists an admissible curve $\gamma$ from $p$ to $q$.
\end{proposition}
\begin{proof}
	Let $p,q\in M$. Since a connected manifold is path-connected, there exists a continuous path $c: [a, b] \to M$ from $p$ to $q$. By compactness, there is a partition of $[a, b]$ such that $c([a_{i-1}, a_i])$ is contained in a single smooth coordinate ball (or half-ball in case of a boundary point) for each $i$. Then we may replace each such curve segment by a straight-line path in coordinates, yielding an admissible curve from $p$ to $q$.
\end{proof}
For convenience, we can define distance on each connected component of a disconnected Riemannian manifold.

\begin{proposition}{Isometry Invariance of the Riemannian Distance Function}{Isometry Invariance of the Riemannian Distance Function}
	Suppose $(M, g)$ and $(\tilde{M}, \tilde{g})$ are connected Riemannian manifolds with or without boundary, and $\varphi: M \to \tilde{M}$ is an isometry. Then for any $x, y \in M$, we have
	\begin{equation}
		d_g(x, y) = d_{\tilde{g}}(\varphi(x), \varphi(y))
	\end{equation}
\end{proposition}
\begin{remark}
	Note that unlike lengths of curves, Riemannian distances are not necessarily preserved by local isometries. For example, take $\mathbb{R} \to S^1$ with the standard metrics.
\end{remark}

We wish to show that the Riemannian distance function turns $M$ into a metric space, whose metric topology is the same as its original manifold topology.

SORRY

\section{Pseudo-Riemannian Metrics}
There is a generalization of Riemannian metrics that has become especially important because of its application to physics.

We use the term scalar product to denote any nondegenerate symmetric bilinear form on a finite-dimensional vector space $V$. A scalar product space is a finite-dimensional vector space endowed with a scalar product. We shall also reuse the notation $\langle \cdot, \cdot \rangle$ for a scalar product, and a norm $|\cdot|$ on a scalar product space is defined by $|v| = \sqrt{|\langle v, v \rangle|}$. Note that in the indefinite case, it is possible for a nonzero vector to be orthogonal to itself, and thus to have norm zero. Thus the norm is not a norm in the usual sense, but we shall still use the term anyway.

If $S \subseteq V$ is any linear subspace, the set of vectors orthogonal to $S$ is denoted by $S^\perp = \{v \in V : \langle v, w \rangle = 0, \forall w \in S\}$. For any finite-dimensional scalar product space, we have
\begin{equation}
	\dim(S) + \dim(S^\perp) = \dim(V), \qquad (S^\perp)^\perp = S
\end{equation}
\begin{remark}
	But usually we cannot decompose $V$ as a direct sum $V = S \oplus S^\perp$ in the indefinite case, because $S \cap S^\perp$ may be nontrivial.
\end{remark}

We wish to prove that every scalar product space has an orthonormal basis. Note that the usual Gram–Schmidt algorithm does not always work in this situation, because the vectors that appear in the denominators might have vanishing norms. We omit the discussion here to leave it to basic linear algebra. The following proposition is a consequence of the existence of orthonormal bases.

\begin{proposition}{Sylvester's Law of Inertia}{Sylvester's Law of Inertia}
	Suppose $(V, q)$ is a finite-dimensional scalar product space. Then there exists a basis $(\beta^1, \ldots, \beta^n)$ for $V^*$ such that
	\begin{equation}
		q = (\beta^1)^2 + \cdots + (\beta^r)^2 - (\beta^{r+1})^2 - \cdots - (\beta^{r+s})^2
	\end{equation}
	for some nonnegative integers $r, s$ with $r + s = \dim(V)$. And the numbers $r$ and $s$ are independent of the choice of basis $(\beta^1, \ldots, \beta^n)$. The pair $(r, s)$ is called the \textbf{signature} of $q$, and $s$ is called the \textbf{index} of $q$.
\end{proposition}

\begin{definition}{Pseudo-Riemannian Manifold}{Pseudo-Riemannian Manifold}
	Suppose $M$ is a smooth manifold. A \textbf{pseudo-Riemannian metric} on $M$ is a smooth symmetric 2-tensor field $g$ on $M$ that is nondegenerate at each point $M$, and with the same signature at each point.
\end{definition}
Every Riemannian metric is a pseudo-Riemannian metric of signature $(n, 0)$.

\begin{proposition}{Existence of Orthonormal Frames}{Existence of Orthonormal Frames}
	Suppose $(M, g)$ is a pseudo-Riemannian manifold. For each point $p \in M$, there exists a smooth orthonormal local frame on a neighborhood of $p$.
\end{proposition}

\begin{proposition}{Product of Pseudo-Riemannian Manifolds}{Product of Pseudo-Riemannian Manifolds}
	If $(M_1, g_1)$ and $(M_2, g_2)$ are pseudo-Riemannian manifolds, then the product manifold $M_1 \times M_2$ has a natural pseudo-Riemannian metric $g = g_1 \oplus g_2$ defined by
	\begin{equation}
		g_{(p, q)}((v_1, v_2), (w_1, w_2)) = g_{1, p}(v_1, w_1) + g_{2, q}(v_2, w_2)
	\end{equation}
	for all $(p, q) \in M_1 \times M_2$, $v_i, w_i \in T_pM_i$.
\end{proposition}

For nonnegative integers $r,s$, we define the \textbf{pseudo-Euclidean space} $\mathbb{R}^{r, s}$ to be the manifold $\mathbb{R}^{r+s}$ with the pseudo-Riemannian metric on standard coordinates $(\xi^1, \ldots, \xi^r, \tau^1, \ldots, \tau^s)$ given by
\begin{equation}
	\bar{q}^{(r, s)} = (\mathrm{d}\xi^1)^2 + \cdots + (\mathrm{d}\xi^r)^2 - (\mathrm{d}\tau^1)^2 - \cdots - (\mathrm{d}\tau^s)^2
\end{equation}

By far the most important pseudo-Riemannian metrics (other than the Riemannian ones) are the \textbf{Lorentz metrics}, which are pseudo-Riemannian metrics of index $1$, and thus the signature is $(r, 1)$.

The pseudo-Euclidean metric $\bar{q}^{(r, 1)}$ is a Lorentz metric called the \textbf{Minkowski metric}, and the pseudo-Euclidean space $\mathbb{R}^{r, 1}$ is called the \textbf{Minkowski space}.
\begin{equation}
	\bar{q}^{(r, 1)} = (\mathrm{d}\xi^1)^2 + \cdots + (\mathrm{d}\xi^r)^2 - (\mathrm{d}\tau)^2
\end{equation}
The Minkowski space $\mathbb{R}^{3, 1}$ is the mathematical model of spacetime in special relativity. The general theory of relativity includes gravitational effects by allowing the Lorentz metric to vary from point to point.

\begin{remark}
	Many, but not all, results from the theory of Riemannian metrics apply equally well to pseudo-Riemannian metrics. Throughout this book, we will attempt to point out which results carry over directly to pseudo-Riemannian metrics, which ones can be adapted with minor modifications, and which ones do not carry over at all. As a rule of thumb, proofs that depend only on the nondegeneracy of the metric tensor, such as properties of the musical isomorphisms and existence and uniqueness of geodesics, work fine in the pseudo-Riemannian setting, while proofs that use positivity in an essential way, such as those involving lengths of curves, do not.
\end{remark}

One notable result that does not carry over to the pseudo-Riemannian case is proposition \ref{prop:Existence of Riemannian Metrics}, which states that every smooth manifold admits a Riemannian metric.

For example, the following result characterizes those manifolds that admit Lorentz metrics.

\begin{theorem}{Existence of Lorentz Metrics}{Existence of Lorentz Metrics}
	A smooth manifold $M$ admits a Lorentz metric if and only if it admits a rank-1 tangent distribution (i.e., a rank-1 subbundle of $TM$).
\end{theorem}
\begin{remark}
	For example, $S^2$ does not admit a Lorentz metric, because it does not admit a rank-1 tangent distribution, because of the hairy ball theorem.

	With some more sophisticated tools from algebraic topology, it can be shown that every noncompact connected smooth manifold admits a Lorentz metric, and a compact connected smooth manifold admits a Lorentz metric if and only if its Euler characteristic is zero. It follows that no even-dimensional sphere admits a Lorentz metric, because its Euler characteristic is $2$.
\end{remark}

\subsection{Pseudo-Riemannian Submanifolds}
The theory of submanifolds is only slightly more complicated in the pseudoRiemannian case. If $(\tilde{M}, \tilde{g})$ is a pseudo-Riemannian manifold, a smooth submanifold $\iota: M \hookrightarrow \tilde{M}$ is called a \textbf{pseudo-Riemannian submanifold} if the pullback $\iota^* \tilde{g}$ is nondegenerate with constant signature. If this is the case, we always consider $M$ to be endowed with the induced pseudo-Riemannian metric $\iota^* \tilde{g}$.
\begin{remark}
	The nondegeneracy hypothesis is not automatically satisfied: Take the $\xi = \tau$ line in $\mathbb{R}^{1, 1}$ with the Minkowski metric. The induced metric on this line is degenerate.
\end{remark}

For hypersurfaces (submanifolds of codimension 1), the nondegeneracy condition is easy to check. If $M \subseteq \tilde{M}$ is a smooth submanifold and $p \in M$, then a vector $v \in T_p \tilde{M}$ is said to be \textbf{normal} to $M$ at $p$ if $\langle v, x \rangle_{\tilde{g}} = 0$ for all $x \in T_pM$. The space of all normal vectors at $p$ is called the \textbf{normal space} to $M$ at $p$, and is denoted by $N_pM$.

\begin{proposition}{hypersurfaces as Pseudo-Riemannian Submanifolds}{hypersurfaces as Pseudo-Riemannian Submanifolds}
	Suppose $(\tilde{M}, \tilde{g})$ is a pseudo-Riemannian manifold of signature $(r, s)$, and $M \subseteq \tilde{M}$ is a smooth hypersurface. Let $\iota: M \hookrightarrow \tilde{M}$ be the inclusion map. Then the pullback $\iota^* \tilde{g}$ is nondegenerate if and only if $\forall p \in M, \forall v \in N_pM \setminus \{0\}, \langle v, v \rangle_{\tilde{g}} \neq 0$.

	If $\tilde{g}(v, v) > 0$ for all $p \in M$ and $v \in N_pM \setminus \{0\}$, then $\iota^* \tilde{g}$ has signature $(r-1, s)$, and if $\tilde{g}(v, v) < 0$ for all $p \in M$ and $v \in N_pM \setminus \{0\}$, then $\iota^* \tilde{g}$ has signature $(r, s-1)$.
\end{proposition}

The gradient, divergence, and Laplacian operators can be defined in the pseudo-Riemannian setting in exactly the same way as in the Riemannian case.

\begin{proposition}{Gradient in Pseudo-Riemannian Manifolds}{Gradient in Pseudo-Riemannian Manifolds}
	Suppose $(\tilde{M}, \tilde{g})$ is a pseudo-Riemannian manifold of signature $(r, s)$, and $f \in C^\infty(\tilde{M})$, and $M = f^{-1}(c)$ for some $c \in \mathbb{R}$. If $\tilde{g}(\grad f, \grad f) > 0$ everywhere on $M$, then $M$ is an embedded pseudo-Riemannian submanifold of $(\tilde{M}, \tilde{g})$ of signature $(r-1, s)$, and if $\tilde{g}(\grad f, \grad f) < 0$ everywhere on $M$, then $M$ is an embedded pseudo-Riemannian submanifold of $(\tilde{M}, \tilde{g})$ of signature $(r, s-1)$. In either case $\grad f$ is a normal vector field on $M$.
\end{proposition}

\begin{proposition}{Pseudo-Riemannian Adapted Orthonormal Frames}{Pseudo-Riemannian Adapted Orthonormal Frames}
	Suppose $(\tilde{M}, \tilde{g})$ is a pseudo-Riemannian manifold and $M \subseteq \tilde{M}$ is an embedded pseudo-Riemannian or Riemannian submanifold. For each point $p \in M$, there exists a smooth orthonormal local frame $(E_1, \ldots, E_n)$ on a neighborhood of $p$ in $\tilde{M}$ that is adapted to $M$.
\end{proposition}

\begin{proposition}{Normal Bundle to a Pseudo-Riemannian Submanifold}{Normal Bundle to a Pseudo-Riemannian Submanifold}
	Suppose $(\tilde{M}, \tilde{g})$ is a pseudo-Riemannian manifold and $M \subseteq \tilde{M}$ is an embedded pseudo-Riemannian or Riemannian submanifold. Then the set of vectors normal to $M$ at each point $p \in M$ forms a smooth vector subbundle of $T\tilde{M}|_M$, called the \textbf{normal bundle} of $M$ in $\tilde{M}$.
\end{proposition}

\section{Other Generalizations of Riemannian Metrics}
Pseudo-Riemannian metrics are obtained by relaxing the positivity requirement for Riemannian metrics. In addition, there are other useful generalizations that result when we relax other requirements.

\subsection{Sub-Riemannian Metrics}
The first generalization arises when we relax the requirement that the metric be defined on the whole tangent space.

\begin{definition}{Sub-Riemannian Metric}{Sub-Riemannian Metric}
	Suppose $M$ is a smooth manifold. A \textbf{sub-Riemannian metric} on $M$ is a smooth fiber metric on a smooth tangent distribution $S \subseteq TM$ (i.e., a vector subbundle of $TM$).
\end{definition}
Since lengths make sense only for vectors in $S$, the only curves whose lengths can be measured are those whose velocity vectors lie everywhere in $S$. Therefore, one usually imposes some condition on $S$ that guarantees that any two nearby points can be connected by such a curve. This is, in a sense, the opposite of the Frobenius integrability condition, which would restrict every such curve to lie in a single leaf of a foliation.

\subsection{Finsler Metrics}
Another important generalization arises from relaxing the requirement that norms of vectors be defined in terms of an inner product on each tangent space.

Suppose $M$ is a smooth manifold. A \textbf{Finsler metric} on $M$ is a continuous function $F: TM \to \mathbb{R}$, smooth on the set of nonzero vectors, whose restriction to each tangent space $T_pM$ is a norm on $T_pM$.

\end{document}
